\documentclass[11pt,reqno]{amsart}

\usepackage[margin=1.08in]{geometry}
\usepackage{amsmath,amssymb,amsthm,mathtools}
\usepackage{booktabs,array,microtype}
\usepackage{xcolor}
\usepackage[colorlinks=true,linkcolor=blue!55!black,citecolor=blue!55!black,
            urlcolor=blue!60!black]{hyperref}

\newtheorem{theorem}{Theorem}[section]
\newtheorem{proposition}[theorem]{Proposition}
\newtheorem{lemma}[theorem]{Lemma}
\newtheorem{corollary}[theorem]{Corollary}
\theoremstyle{definition}
\newtheorem{definition}[theorem]{Definition}
\theoremstyle{remark}
\newtheorem{remark}[theorem]{Remark}
\newtheorem{status}[theorem]{Status}

\newcommand{\Q}{\mathbb Q}
\newcommand{\Z}{\mathbb Z}
\newcommand{\Zp}{\mathbb Z_{(p)}}
\newcommand{\bin}[2]{\binom{#1}{#2}}
\newcommand{\Hd}[2]{H_{#1}^{(#2)}}
\newcommand{\Hh}[1]{H_{#1}}
\newcommand{\Ad}[2]{A_{#1}(#2)}
\newcommand{\Bd}[2]{B_{#1}(#2)}
\DeclareMathOperator{\ord}{ord}

\title[Harmonic companion rows]{Harmonic companion rows for binomial and
Ap\'ery-like sums:\\closed forms, Lucas descent, and the $\zeta(2)$ boundary}
\author{}
\date{August 2026}

\begin{document}
\begin{abstract}
For $d\geq1$ let $A_d(n)=\sum_k\binom nk^d$.  We prove a Lucas theorem for
every even coefficient of the Straub--Zudilin deformation of $A_d$.  In
particular, its normalized quadratic coefficient is
\[
 B_d(n)=\sum_{k=0}^n\binom nk^d\left\{
 \frac{H_k^{(2)}+H_{n-k}^{(2)}}{2(d+1)}
 +\frac{d(H_k-H_{n-k})^2}{2(d+1)}\right\}.
\]
and satisfies the uniform digit law
\[
 p^2B_d(ap+r)\equiv B_d(a)A_d(r)\pmod p
 \qquad(0\leq a,r<p,\ p\nmid2(d+1)).
\]
More generally the coefficient of $t^{2j}$ satisfies a $p^{2j}$ digit law,
with an explicit weight given by a complete Bell polynomial in harmonic
numbers.  For $d=3$ and $4$, finite certificates identify the quadratic rows with the
canonical second solutions of the Franel and $s_{10}$ recurrences.  Both
identifications and their recurrence-level Lucas laws have also been checked in
Lean, including every boundary cell.  We explain the rank-two harmonic
certificate mechanism and derive exact Casoratian series, alternating error
bounds, and a polynomial continued fraction for $5/\zeta(2)$.  We compare the
result with the minimal weight-three formula for
Ap\'ery's $\zeta(3)$ companion, and settle the previously open minimal
$\zeta(2)$ formula.  We also obtain a new five-term harmonic companion for
Cooper's sporadic sequence $s_7$.  It proves a $p^2$ companion Lucas law, the
Ap\'ery limit $\zeta(2)/7$, and an associated polynomial continued fraction.
The minimal Ap\'ery proof lifts the harmonic weight to a proper
two-parameter hypergeometric deformation, differentiates an order-three
creative telescope, and eliminates the auxiliary derivatives in the Ore
algebra of shifts.
\end{abstract}
\maketitle

\section{Introduction and main results}

Put
\[
 \Hd m q=\sum_{j=1}^m\frac1{j^q},\qquad \Hd 0q=0,
\]
and write $H_m=H_m^{(1)}$.  The sums
\[
 \Ad dn=\sum_{k=0}^n\bin nk^d
\]
include the Franel numbers at $d=3$ and the sporadic sequence usually denoted
$s_{10}$ at $d=4$.  Straub and Zudilin~\cite{StraubZudilin} embedded these rows
in an analytic deformation and determined its Ap\'ery limits.  Its normalized
quadratic coefficient is the following entirely finite sum.

\begin{definition}[harmonic companion row]\label{def:Bd}
For $d\geq1$, set
\begin{equation}\label{eq:Bd}
 \Bd dn=\sum_{k=0}^n\bin nk^d w_d(n,k),\qquad
 w_d(n,k)=\frac{\sigma(n,k)+d\delta(n,k)^2}{2(d+1)},
\end{equation}
where
\begin{equation}\label{eq:sigdel}
 \sigma(n,k)=\Hd k2+\Hd{n-k}2,
 \qquad \delta(n,k)=H_k-H_{n-k}.
\end{equation}
Thus $B_d(0)=0$ and $B_d(1)=1$.
\end{definition}

Our first result does not require a recurrence.

\begin{theorem}[uniform Lucas law]\label{thm:universal}
Let $d\geq1$, let $p$ be a prime with $p\nmid2(d+1)$, and let
$0\leq a,r<p$.  Then $p^2B_d(ap+r)$ is $p$-integral and
\begin{equation}\label{eq:universal}
 p^2\Bd d{ap+r}\equiv \Bd da\Ad dr\pmod {p\Zp}.
\end{equation}
\end{theorem}

Here and below a rational congruence $x\equiv y\pmod{p\Zp}$ means that both
numbers are in the localization $\Zp$ and $x-y\in p\Zp$.  This convention is
essential because harmonic numbers need not themselves be $p$-integral.

For $d=3,4$ the companion rows are genuine second solutions of order-two
recurrences.

\begin{theorem}[Franel companion]\label{thm:franel}
Let $F_n=\sum_k\bin nk^3$ and let $C_0=0,C_1=1$ satisfy
\begin{equation}\label{eq:Frec}
 (n+1)^2C_{n+1}=(7n^2+7n+2)C_n+8n^2C_{n-1}\qquad(n\geq1).
\end{equation}
Then
\begin{equation}\label{eq:Fform}
 C_n=\sum_{k=0}^n\bin nk^3
 \left\{\frac{\Hd k2+\Hd{n-k}2}{8}
       +\frac38(H_k-H_{n-k})^2\right\}=B_3(n).
\end{equation}
For every odd prime $p$ and $0\leq a,r<p$,
\begin{equation}\label{eq:Flucas}
 p^2C_{ap+r}\equiv C_aF_r\pmod {p\Zp}.
\end{equation}
\end{theorem}

\begin{theorem}[$s_{10}$ companion]\label{thm:s10}
Let $S_n=\sum_k\bin nk^4$ and let $D_0=0,D_1=1$ satisfy
\begin{equation}\label{eq:s10rec}
 (n+1)^3D_{n+1}=(2n+1)(6n^2+6n+2)D_n
                 +n(64n^2-4)D_{n-1}.
\end{equation}
Then
\begin{equation}\label{eq:s10form}
 D_n=\sum_{k=0}^n\bin nk^4
 \left\{\frac{\Hd k2+\Hd{n-k}2}{10}
       +\frac25(H_k-H_{n-k})^2\right\}=B_4(n).
\end{equation}
If $p\nmid10$ and $0\leq a,r<p$, then
\begin{equation}\label{eq:s10lucas}
 p^2D_{ap+r}\equiv D_aS_r\pmod {p\Zp}.
\end{equation}
\end{theorem}

The harmonic formulas in Theorems~\ref{thm:franel} and~\ref{thm:s10} are the
$d=3,4$ specializations of the known deformation coefficient.  The additions
here are (i) the direct Lucas descent for the normalized companion row,
(ii) explicit finite all-boundary recurrence certificates in these two cases,
and (iii) a machine-checked identification with the recurrence-defined second
solutions.

\section{The deformation and the universal weight}

Consider the Straub--Zudilin deformation
\begin{equation}\label{eq:deformation}
 \mathcal A_d(n,t)=\sum_{k=0}^n\bin nk^d
 \prod_{j=1}^k\left(1-\frac tj\right)^{-d}
 \prod_{j=1}^{n-k}\left(1+\frac tj\right)^{-d}.
\end{equation}
The substitution $k\mapsto n-k$, $t\mapsto-t$ shows that the sum is even in
$t$.  For a fixed summand, logarithmic expansion gives
\[
 \begin{split}
 &\prod_{j=1}^k(1-t/j)^{-d}
  \prod_{j=1}^{n-k}(1+t/j)^{-d}\\
 &\qquad=\exp\left(dt\delta+\frac d2t^2\sigma+O(t^3)\right)
 =1+dt\delta+\frac{t^2}{2}\bigl(d\sigma+d^2\delta^2\bigr)+O(t^3).
 \end{split}
\]
Consequently, if $\mathcal A_{d,1}(n)$ denotes the coefficient of $t^2$ in
\eqref{eq:deformation}, then
\begin{equation}\label{eq:coefficient}
 \mathcal A_{d,1}(n)=d(d+1)B_d(n).
\end{equation}
The normalization is forced by $\mathcal A_{d,1}(1)=d(d+1)$.

Straub and Zudilin proved that deformation coefficients solve the telescoping
recurrences attached to $A_d$ in their stated stable range and that
\begin{equation}\label{eq:limit}
 \lim_{n\to\infty}\frac{B_d(n)}{A_d(n)}=\frac{\zeta(2)}{d+1}.
\end{equation}
The limit is also transparent directly from~\eqref{eq:Bd}: the binomial-power
measure concentrates at $k=n/2$, where $\sigma(n,k)\to2\zeta(2)$ and
$\delta(n,k)\to0$.  Formula~\eqref{eq:limit} is contextual rather than a new
claim here.

For $d\geq5$ the minimal recurrence for $A_d$ need not have order two, so the
phrase ``second solution'' is no longer canonical.  Definition~\ref{def:Bd}
and Theorem~\ref{thm:universal}, however, remain meaningful for every $d$.

\section{Proof of the universal Lucas law}

The proof uses only ordinary Lucas congruence and harmonic digit descent.

\begin{lemma}[harmonic descent]\label{lem:harmonic}
For a prime $p$, $q\geq1$, $b\geq0$, and $0\leq s<p$,
\begin{equation}\label{eq:hdesc}
 p^q H_{bp+s}^{(q)}\in\Zp,
 \qquad p^q H_{bp+s}^{(q)}\equiv H_b^{(q)}\pmod {p\Zp}.
\end{equation}
\end{lemma}

\begin{proof}
Separate the summation indices divisible by $p$:
\[
 p^q\sum_{j=1}^{bp+s}\frac1{j^q}
 =\sum_{i=1}^b\frac1{i^q}
  +\sum_{\substack{1\leq j\leq bp+s\\p\nmid j}}\frac{p^q}{j^q}.
\]
Every term in the second sum belongs to $p\Zp$.
\end{proof}

\begin{proof}[Proof of Theorem~\ref{thm:universal}]
Write every $k\leq ap+r$ uniquely as $k=bp+s$, with $0\leq s<p$.
By Lemma~\ref{lem:harmonic}, $p^2w_d(ap+r,k)$ is $p$-integral.  If $s>r$,
Lucas' theorem gives $p\mid\binom{ap+r}{bp+s}$; hence
\[
 \binom{ap+r}{bp+s}^{d}p^2w_d(ap+r,bp+s)\equiv0\pmod {p\Zp}.
\]
If $s\leq r$, subtraction has no borrow:
\[
 ap+r-(bp+s)=(a-b)p+(r-s).
\]
Lucas' theorem and Lemma~\ref{lem:harmonic} now give
\begin{align*}
 \binom{ap+r}{bp+s}^{d}&\equiv\binom ab^d\binom rs^d\pmod p,\\
 p^2\sigma(ap+r,bp+s)&\equiv H_b^{(2)}+H_{a-b}^{(2)}\pmod {p\Zp},\\
 p^2\delta(ap+r,bp+s)^2&\equiv(H_b-H_{a-b})^2\pmod {p\Zp}.
\end{align*}
Because $2(d+1)$ is a unit in $\Zp$, the last two lines say
$p^2w_d(ap+r,bp+s)\equiv w_d(a,b)$.  Summing the surviving digit rectangle
$0\leq b\leq a$, $0\leq s\leq r$ factors the result:
\[
 p^2B_d(ap+r)\equiv
 \sum_{b=0}^a\sum_{s=0}^r\binom ab^d\binom rs^d w_d(a,b)
 =B_d(a)A_d(r).
\]
\end{proof}

\begin{remark}
The proof explains both the exponent $p^2$ and the shape of the answer.  Each
monomial in the weight is homogeneous of harmonic weight two, so multiplication
by $p^2$ extracts its high digit.  Cells with a low-digit borrow disappear
through the binomial shell; the remaining rectangle separates into high and low
digit sums.
\end{remark}

\section{The full higher-coefficient Lucas tower}\label{sec:tower}

The quadratic row is the first case of a theorem for every coefficient of the
deformation.  Put
\begin{equation}\label{eq:Lq}
 \mathcal L_q(n,k)=H_k^{(q)}+(-1)^qH_{n-k}^{(q)}
\end{equation}
and define the homogeneous polynomial
\begin{equation}\label{eq:Qdm}
 \mathcal Q_{d,m}(n,k)
 =[t^m]\exp\left(d\sum_{q\geq1}\frac{\mathcal L_q(n,k)}q,t^q\right).
\end{equation}
Equivalently,
\begin{equation}\label{eq:partitionQ}
 \mathcal Q_{d,m}
 =\sum_{m_1+2m_2+\cdots+mm_m=m}
   \prod_{q=1}^{m}\frac1{m_q!}
   \left(\frac{d\mathcal L_q}{q}\right)^{m_q}.
\end{equation}
Every monomial on the right has total harmonic weight $m$.  By logarithmic
expansion of~\eqref{eq:deformation}, its coefficient rows are
\begin{equation}\label{eq:Adj}
 A_{d,j}^{\langle2\rangle}(n)
 :=[t^{2j}]\mathcal A_d(n,t)
 =\sum_{k=0}^n\binom nk^d\mathcal Q_{d,2j}(n,k).
\end{equation}

\begin{theorem}[higher-coefficient Lucas tower]\label{thm:tower}
Let $d,j\geq1$, let $p>2j$ be prime, and let $0\leq a,r<p$.  Then
\begin{equation}\label{eq:tower}
 p^{2j}A_{d,j}^{\langle2\rangle}(ap+r)
 \equiv A_{d,j}^{\langle2\rangle}(a)A_d(r)\pmod {p\Zp}.
\end{equation}
If
\begin{equation}\label{eq:normalizer}
 c_{d,j}:=A_{d,j}^{\langle2\rangle}(1)
 =2\binom{d+2j-1}{2j}
\end{equation}
is a $p$-adic unit, the normalized row
$\widehat A_{d,j}=A_{d,j}^{\langle2\rangle}/c_{d,j}$ satisfies
\begin{equation}\label{eq:towernorm}
 p^{2j}\widehat A_{d,j}(ap+r)
 \equiv\widehat A_{d,j}(a)A_d(r)\pmod {p\Zp}.
\end{equation}
\end{theorem}

\begin{proof}
For $q\leq2j$, Lemma~\ref{lem:harmonic} gives
\[
 p^q\mathcal L_q(ap+r,bp+s)
 \equiv\mathcal L_q(a,b)\pmod {p\Zp}
\]
whenever $s\leq r$, since then
$ap+r-(bp+s)=(a-b)p+(r-s)$.  Formula~\eqref{eq:partitionQ} is homogeneous of
weight $2j$; its denominators $q$ and $m_q!$ are $p$-adic units because
$p>2j$.  Hence
\begin{equation}\label{eq:Qdescent}
 p^{2j}\mathcal Q_{d,2j}(ap+r,bp+s)
 \equiv\mathcal Q_{d,2j}(a,b)\pmod {p\Zp}.
\end{equation}
For a borrow cell $s>r$, $p^{2j}\mathcal Q_{d,2j}$ is $p$-integral and
Lucas' theorem makes $\binom{ap+r}{bp+s}^d$ divisible by $p$.  Those cells
therefore vanish.  On the remaining rectangle Lucas' theorem and
\eqref{eq:Qdescent} give
\begin{align*}
 p^{2j}A_{d,j}^{\langle2\rangle}(ap+r)
 &\equiv\sum_{b=0}^a\sum_{s=0}^r
   \binom ab^d\binom rs^d\mathcal Q_{d,2j}(a,b)\\
 &=A_{d,j}^{\langle2\rangle}(a)A_d(r).
\end{align*}
At $n=1$, equation~\eqref{eq:deformation} is
$(1-t)^{-d}+(1+t)^{-d}$, whose $t^{2j}$ coefficient is
$2\binom{d+2j-1}{2j}$.  Division by a $p$-adic unit proves the normalized
statement.
\end{proof}

The coefficientwise tower has a useful generating-series form.  Adopt the
convention $A_{d,0}^{\langle2\rangle}=A_d$ and, for a prime $p$, put
\begin{equation}\label{eq:truncatedFrob}
 \mathcal F_{d,p}^{[J]}(n;T)
 =\sum_{j=0}^{J}p^{2j}A_{d,j}^{\langle2\rangle}(n)T^{2j}.
\end{equation}

\begin{corollary}[truncated generating-series Lucas law]\label{cor:seriesLucas}
Let $d,J\geq1$, let $p>2J$ be prime, and let $0\leq a,r<p$.  Then, coefficientwise
in $(\Zp/p\Zp)[T]$,
\begin{equation}\label{eq:seriesLucas}
 \boxed{\displaystyle
 \mathcal F_{d,p}^{[J]}(ap+r;T)
 \equiv A_d(r)\mathcal F_{d,p}^{[J]}(a;T)\pmod p.}
\end{equation}
Thus the rescaled truncated deformation, rather than any single harmonic row,
is itself a Lucas object.
\end{corollary}

\begin{proof}
The constant coefficient is the ordinary Lucas congruence for $A_d$.  For
$1\leq j\leq J$, the coefficient of $T^{2j}$ is exactly
Theorem~\ref{thm:tower}; the common hypothesis $p>2J$ implies $p>2j$.
\end{proof}

Theorem~\ref{thm:universal} is the case $j=1$, because
\[
 \mathcal Q_{d,2}=\frac d2\sigma+\frac{d^2}{2}\delta^2
 =d(d+1)w_d.
\]
The first genuinely new weight occurs at $j=2$.  If
$\delta_q=H_k^{(q)}-H_{n-k}^{(q)}$ and
$\sigma_q=H_k^{(q)}+H_{n-k}^{(q)}$, then
\begin{equation}\label{eq:Q4}
 \mathcal Q_{d,4}
 =\frac d4\sigma_4+\frac{d^2}{3}\delta_1\delta_3
  +\frac{d^2}{8}\sigma_2^2+\frac{d^3}{4}\delta_1^2\sigma_2
  +\frac{d^4}{24}\delta_1^4.
\end{equation}

\begin{corollary}[the $p^4$ row]\label{cor:p4}
Let
\[
 \widehat A_{d,2}(n)=
 \frac{12}{d(d+1)(d+2)(d+3)}
 \sum_{k=0}^n\binom nk^d\mathcal Q_{d,4}(n,k).
\]
For $p>4$, $0\leq a,r<p$, and
$p\nmid d(d+1)(d+2)(d+3)$,
\[
 p^4\widehat A_{d,2}(ap+r)
 \equiv\widehat A_{d,2}(a)A_d(r)\pmod {p\Zp}.
\]
Moreover, after Straub--Zudilin~\cite{StraubZudilin},
\begin{equation}\label{eq:zeta4limit}
 \lim_{n\to\infty}\frac{\widehat A_{d,2}(n)}{A_d(n)}
 =\frac{3(5d+2)}{(d+1)(d+2)(d+3)}\zeta(4).
\end{equation}
For $d\geq5$ the row solves every telescoping recurrence in the stable range
specified in their theorem.
\end{corollary}

\begin{proof}
The congruence is Theorem~\ref{thm:tower}.  The coefficient of $t^4$ in
$(t/\sin t)^d$ is $d(5d+2)/360$; divide the resulting limit
$d(5d+2)\pi^4/360$ by
$c_{d,2}=d(d+1)(d+2)(d+3)/12$ and use $\zeta(4)=\pi^4/90$.
\end{proof}

\section{Finite recurrence certificates for \texorpdfstring{$d=3,4$}{d=3,4}}

The analytic deformation proves recurrence membership in a stable range.  For
the Franel and $s_{10}$ rows we instead use a finite identity valid at every
$n\geq1$, including the top cells.  We summarize the method because it also
explains why the harmonic weight in~\eqref{eq:Bd} is tractable.

Let $S_d(n,k)=\binom nk^d$ and let
\[
 L[u]_n=c_+(n)u_{n+1}+c_0(n)u_n+c_-(n)u_{n-1}.
\]
A Zeilberger certificate $G$ is chosen so that
\begin{equation}\label{eq:Zcert}
 L[S_d(\cdot,k)]_n=G(n,k+1)-G(n,k),
 \qquad
 G(n,k)=S_d(n,k)\frac{k^d\rho(n,k)}{(n+1-k)^d}.
\end{equation}
The symmetric letters in~\eqref{eq:sigdel} obey
\begin{align}
 \sigma(n+1,k)-\sigma(n,k)&=(n+1-k)^{-2},&
 \delta(n+1,k)-\delta(n,k)&=-(n+1-k)^{-1},\label{eq:ndiff}\\
 \sigma(n,k+1)-\sigma(n,k)&=(k+1)^{-2}-(n-k)^{-2},&
 \delta(n,k+1)-\delta(n,k)&=(k+1)^{-1}+(n-k)^{-1}.\label{eq:kdiff}
\end{align}
Thus, after Abel summation, every residual term belongs to the rank-two module
\[
 S_d(n,k)\bigl(\Q(n,k)+\Q(n,k)\delta(n,k)\bigr).
\]
It is enough to solve two rational Gosper equations and form
\begin{equation}\label{eq:Psi}
 \Psi(n,k)=S_d(n,k)\bigl(z_0(n,k)+z_1(n,k)\delta(n,k)\bigr).
\end{equation}
The remaining assertion is a finite list of rational identities plus boundary
values at $k=0,n,n+1$.  Clearing denominators turns each one into a polynomial
identity over $\Q$.

For reference, the Franel Zeilberger numerator is
\begin{align}\label{eq:rhoF}
 \rho_F(n,k)=-\frac1n\bigl(&4-12k+12k^2-4k^3+22n-39kn+18k^2n\\
 &+32n^2-27kn^2+14n^3\bigr).
\end{align}
The two Gosper numerators may be taken as
\begin{align*}
N_F=\frac{3k^2}{4n}\bigl(&4-16k+24k^2-16k^3+4k^4+26n-68kn+63k^2n-20k^3n\\
 &+54n^2-88kn^2+39k^2n^2+46n^3-36kn^3+14n^4\bigr),\\
M_F=-\frac{k}{2n}\bigl(&2-10k+20k^2-20k^3+10k^4-2k^5+15n-50kn+70k^2n\\
 &-45k^3n+11k^4n+40n^2-90kn^2+80k^2n^2-25k^3n^2\\
 &+50n^3-70kn^3+30k^2n^3+30n^4-20kn^4+7n^5\bigr),
\end{align*}
with $z_1=N_F/(n+1-k)^4$ and $z_0=M_F/(n+1-k)^5$.
They satisfy the forced top values
\[
 \rho_F(n,n+1)=-(n+1)^2,
 \quad N_F(n)=\frac34\bigl((n+1)^5-(n+1)\bigr),
 \quad M_F(n)=-\frac12\bigl((n+1)^5+1\bigr).
\]
The $d=4$ certificate has the same architecture but substantially larger
polynomials; it is recorded exactly in the accompanying Lean source.  Its top
values are
\[
 \rho_{10}(n,n+1)=-(n+1)^3,
 \quad N_{10}(n)=\frac45\bigl((n+1)^7-(n+1)^2\bigr),
 \quad M_{10}(n)=-\frac12\bigl((n+1)^7+(n+1)\bigr).
\]
Equations~\eqref{eq:Zcert}--\eqref{eq:Psi}, the boundary identities, and the
initial values prove Theorems~\ref{thm:franel} and~\ref{thm:s10};
Theorem~\ref{thm:universal} then gives their Lucas congruences.

The first values provide a useful normalization check:
\[
 \begin{array}{c|rrrrrr}
 n&0&1&2&3&4&5\\ \hline
 C_n&0&1&4&208/9&1280/9&208384/225
 \end{array}
\]
and
\[
 \begin{array}{c|rrrrr}
 n&0&1&2&3&4\\ \hline
 D_n&0&1&21/4&1001/18&85085/144.
 \end{array}
\]

\section{A new harmonic companion for Cooper's \texorpdfstring{$s_7$}{s7}}
\label{sec:s7}

The harmonic-lifting method is not confined to powers of a single binomial
coefficient.  Consider Cooper's sporadic sequence
$s_7$~\cite{Gorodetsky,MalikStraub},
\begin{equation}\label{eq:s7shell}
 A_n^{(7)}=\sum_{k=0}^n S_7(n,k),\qquad
 S_7(n,k)=\bin nk^2\bin{n+k}k\bin{2k}n.
\end{equation}
It satisfies
\begin{equation}\label{eq:s7rec}
 (n+1)^3u_{n+1}=(2n+1)(13n^2+13n+4)u_n
                  +3n(9n^2-1)u_{n-1}.
\end{equation}
Let $B_0^{(7)}=0,B_1^{(7)}=1$ denote its normalized companion solution.

\begin{theorem}[the $s_7$ harmonic companion]\label{thm:s7harm}
For every $n\geq0$,
\begin{equation}\label{eq:s7harm}
 \boxed{\displaystyle
 B_n^{(7)}=\frac1{28}\sum_{k=0}^nS_7(n,k)
 \left\{
 3(H_k-H_{n-k})(2H_{2k}-3H_k+H_{n-k})
 +5H_k^{(2)}+2H_n^{(2)}-3H_{n-k}^{(2)}
 \right\}.}
\end{equation}
The proof also gives the auxiliary identity
\begin{equation}\label{eq:s7zero}
 \sum_{k=0}^nS_7(n,k)(2H_{2k}-3H_k+H_{n-k})=0.
\end{equation}
\end{theorem}

\begin{proof}
Put
\[
 X_m(t)=\frac{m!}{(1-t)_m}
       =\prod_{j=1}^m(1-t/j)^{-1}
\]
and lift the summand to the proper hypergeometric deformation
\begin{align}\label{eq:s7Phi}
 \Phi_n(u,v)=\sum_k S_7(n,k)\,&
 \frac{X_k(3u+5v/3)}{X_k(14v/3)}X_{n-k}(-3u+v)X_{2k}(2v)\notag\\
 &\times\frac{X_n(2u+2v)}{X_n(u+v)^2}.
\end{align}
At the origin the two logarithmic derivatives of its summand are
\[
 \mathcal A=3(H_k-H_{n-k}),\qquad
 \mathcal B=2H_{2k}-3H_k+H_{n-k},
\]
and its mixed logarithmic derivative is
\[
 \mathcal C=5H_k^{(2)}+2H_n^{(2)}-3H_{n-k}^{(2)}.
\]
Consequently the mixed derivative of the summand is
$S_7(\mathcal A\mathcal B+\mathcal C)$, exactly twenty-eight times the
summand in~\eqref{eq:s7harm}.

It remains to prove recurrence membership; a single direct harmonic Gosper
telescope does not exist.  We instead polarize the mixed derivative.  Set
\[
 \Phi^u_n(t)=\Phi_n(t,0),\qquad
 \Phi^v_n(t)=\Phi_n(0,t),\qquad
 \Phi^d_n(t)=\Phi_n(t,t).
\]
The diagonal is especially small:
\begin{equation}\label{eq:s7diag}
 \Phi^d_n(t)=\sum_kS_7(n,k)
 \frac{X_{n-k}(-2t)X_{2k}(2t)X_n(4t)}{X_n(2t)^2}.
\end{equation}
Exact creative telescoping gives parameter-dependent operators of orders
$2,4,3$ for $\Phi^u,\Phi^v,\Phi^d$, respectively.  Write
$C_r^x=\partial_t^rC^x(t)|_{t=0}$ for their operator coefficients and put
\[
 U=\partial_u\Phi|_0,\quad V=\partial_v\Phi|_0,\quad
 U_2=\partial_u^2\Phi|_0,\quad V_2=\partial_v^2\Phi|_0,
 \quad D_2=\frac{d^2}{dt^2}\Phi(t,t)\big|_{t=0}.
\]
For the forward operator
\begin{equation}\label{eq:s7L}
 \mathcal L_7=-3(n+1)(3n+2)(3n+4)
 -(2n+3)(13n^2+39n+30)E+(n+2)^3E^2,
\end{equation}
exact right division gives
\begin{align*}
 C_0^u&=-7(n+1)^2(n+2)^3\mathcal L_7,\\
 C_0^v&=Q_0\mathcal L_7, & C_1^v&=Q_1\mathcal L_7.
\end{align*}
for explicit order-two rational operators $Q_0,Q_1$.  The first derivative
of the $v$-telescope therefore gives $C_0^vV=0$.  The leading coefficient
does not vanish for $n\geq0$, while $V_0=V_1=V_2=V_3=0$; this proves
\eqref{eq:s7zero}.

The second derivatives of the three telescopes now read
\begin{align*}
 C_0^uU_2+2C_1^uU+C_2^uA^{(7)}&=0,\\
 C_0^vV_2+C_2^vA^{(7)}&=0,\\
 C_0^dD_2+2C_1^dU+C_2^dA^{(7)}&=0.
\end{align*}
Eliminating $A^{(7)}$ and $U$ in the Ore algebra gives
\begin{equation}\label{eq:s7Ore}
 \mathcal P\mathcal L_7J=0,qquad
 J=\frac{D_2-U_2-V_2}{2}=\partial_u\partial_v\Phi|_0,
\end{equation}
where $\mathcal P=P_0+P_1E+P_2E^2$ and
\begin{align*}
 P_0={}&2(n+1)(n+2)(98n^3+1029n^2+3526n+3931),\\
 P_1={}&(n+2)(392n^4+4606n^3+19543n^2+35229n+22598),\\
 P_2={}&(n+3)(2n+7)(98n^3+735n^2+1762n+1336).
\end{align*}
All three cell identities, their Taylor coefficients, and the right remainder
in~\eqref{eq:s7Ore} are exact rational identities.  At integer $n$ the
certificate flux has finite support; equivalently, one may regularize
$n$ to $n+\varepsilon$ and take the removable limits at the exceptional
cells.  Thus summing the cell identities introduces no endpoint term.

The order-four operator $\mathcal P\mathcal L_7$ has nonzero leading
coefficient for every $n\geq0$.  Finally,
\[
 (J_0,J_1,J_2,J_3)=(0,28,315,45010/9)
 =28(B_0^{(7)},B_1^{(7)},B_2^{(7)},B_3^{(7)}).
\]
Since $\mathcal L_7B^{(7)}=0$, four initial values prove $J=28B^{(7)}$.
\end{proof}

\begin{theorem}[$s_7$ companion Lucas law]\label{thm:s7lucas}
Let $p\nmid28$ be prime and let $0\leq a,r<p$.  Then
\begin{equation}\label{eq:s7lucas}
 p^2B_{ap+r}^{(7)}\equiv B_a^{(7)}A_r^{(7)}\pmod {p\Zp}.
\end{equation}
\end{theorem}

\begin{proof}
Write $k=bp+s$, $0\leq s<p$.  The three binomial factors in~\eqref{eq:s7shell}
give the termwise Lucas congruence
\begin{equation}\label{eq:s7termLucas}
 S_7(ap+r,bp+s)\equiv S_7(a,b)S_7(r,s)\pmod p.
\end{equation}
Indeed, if the left side is nonzero modulo $p$, the low digit of
$\binom{ap+r}{bp+s}$ forces $s\leq r$, while that of
$\binom{(a+b)p+r+s}{bp+s}$ forces $r+s<p$.  Hence $2s<p$ as well, and ordinary
Lucas factorization applies without a low-digit carry to all three factors.
If either condition fails, the corresponding low-digit factor and the right
side of~\eqref{eq:s7termLucas} both vanish.

Let $w_7(n,k)$ denote the expression in braces in~\eqref{eq:s7harm}, divided
by $28$.  Lemma~\ref{lem:harmonic} shows that $p^2w_7(ap+r,bp+s)$ is
$p$-integral.  On a nonvanishing Lucas cell we have $s\leq r$ and $2s<p$, so
the same lemma, now for the indices $k,n-k,2k,n$, gives
\[
 p^2w_7(ap+r,bp+s)\equiv w_7(a,b)\pmod {p\Zp}.
\]
All other cells vanish modulo $p$.  Summing the surviving digit rectangle and
using~\eqref{eq:s7termLucas} gives
\[
 p^2B_{ap+r}^{(7)}\equiv
 \sum_{b=0}^a\sum_{s=0}^rS_7(a,b)S_7(r,s)w_7(a,b)
 =B_a^{(7)}A_r^{(7)}.
\]
\end{proof}

\begin{corollary}[the $s_7$ Ap\'ery limit]\label{cor:s7limit}
\begin{equation}\label{eq:s7limit}
 \lim_{n\to\infty}\frac{B_n^{(7)}}{A_n^{(7)}}=\frac{\zeta(2)}7.
\end{equation}
\end{corollary}

\begin{proof}
The factorial form of the shell is
\[
 S_7(n,k)=\frac{(n+k)!(2k)!}{k!^3(n-k)!^2(2k-n)!}.
\]
Its consecutive quotient is
\[
 \frac{S_7(n,k+1)}{S_7(n,k)}
 =\frac{(n+k+1)(2k+2)(2k+1)(n-k)^2}
 {(k+1)^3(2k-n+1)(2k-n+2)}.
\]
For $k/n\to x\in(1/2,1)$ this tends to
$4(1+x)(1-x)^2/[x(2x-1)^2]$, and the equation that this limit equal one
reduces to $4-5x=0$.  The quotient bounds therefore give exponential
concentration at $k/n=4/5$.

At that saddle,
\[
 H_k-H_{n-k}\longrightarrow\log4,qquad
 2H_{2k}-3H_k+H_{n-k}\longrightarrow0,qquad
 H_m^{(2)}\longrightarrow\zeta(2)
\]
for each linearly growing index $m$.  The complement of a fixed saddle
neighborhood is exponentially negligible, even after the $O(\log^2 n)$
harmonic weight is included.  Substitution in~\eqref{eq:s7harm} leaves
$(5+2-3)\zeta(2)/28=\zeta(2)/7$.
\end{proof}

\section{Formal verification}

The Franel and $s_{10}$ recurrence proofs and their two specialized Lucas laws
have been formalized in Lean~4 with Mathlib.  The $s_7$ theorem is presently
verified by the exact symbolic audit described above, but not yet in Lean.
The relevant declarations in
\nolinkurl{lean/ZetaLucas/FranelClosedForm.lean} are
\[
\begin{array}{ll}
\texttt{franel\_harmonic\_closed\_form},&
\texttt{s10\_harmonic\_closed\_form},\\
\texttt{franelB\_lucas},&\texttt{s10B\_lucas}.
\end{array}
\]
The proof does not infer recurrence membership from a finite sample.  Lean
checks the generic cleared-denominator identities, the exceptional boundary
cells, the telescoping finite sums, and the initial-value identification.  The
Lucas statements are transferred to the recurrence-defined sequences only
after the harmonic sums have been identified with those sequences.

The target reproduces with
\begin{verbatim}
cd lean
lake build ZetaLucas.FranelClosedForm
\end{verbatim}
and contains no \texttt{sorry}, \texttt{admit}, custom axiom, or unsafe
declaration.  The axiom audit of the four headline theorems reports only
Mathlib's standard logical axioms \texttt{propext},
\texttt{Classical.choice}, and \texttt{Quot.sound}.

\section{Casoratians, certified brackets, and continued fractions}\label{sec:cas}

The closed forms identify the companion rows, but the recurrence alone then
supplies another exact description.  We record the general observation first.

\begin{lemma}[Casoratian and alternating error]\label{lem:cas}
Suppose positive sequences $A_n$ and $B_n$ satisfy
\begin{equation}\label{eq:genericrec}
 \alpha_nu_{n+1}=\beta_nu_n+\gamma_nu_{n-1}\qquad(n\geq1),
\end{equation}
where $\alpha_n,\beta_n,\gamma_n>0$, and suppose
$A_0=1$, $B_0=0$, $B_1=1$.  Put
\[
 W_n=A_nB_{n+1}-A_{n+1}B_n.
\]
Then
\begin{equation}\label{eq:Wgeneric}
 W_0=1,\qquad W_n=-\frac{\gamma_n}{\alpha_n}W_{n-1},
\end{equation}
and
\begin{equation}\label{eq:increment}
 \frac{B_{n+1}}{A_{n+1}}-\frac{B_n}{A_n}
 =\frac{W_n}{A_nA_{n+1}}.
\end{equation}
If $B_n/A_n\to L$, then
\begin{equation}\label{eq:casSeries}
 L=\sum_{m=0}^{\infty}\frac{W_m}{A_mA_{m+1}}.
\end{equation}
Moreover, with $R_n=B_n/A_n$ and
$\varepsilon_n=|W_n|/(A_nA_{n+1})$,
\begin{equation}\label{eq:brackets}
 \begin{cases}
  R_n<L<R_n+\varepsilon_n,&n\text{ even},\\
  R_n-\varepsilon_n<L<R_n,&n\text{ odd}.
 \end{cases}
\end{equation}
\end{lemma}

\begin{proof}
Substitute the recurrence for $A_{n+1}$ and $B_{n+1}$ into the determinant;
the $\beta_n$ terms cancel and give~\eqref{eq:Wgeneric}.  Division by
$A_nA_{n+1}$ gives~\eqref{eq:increment}, which telescopes.  The terms alternate
by~\eqref{eq:Wgeneric}, and their absolute values strictly decrease because
\[
 \frac{\varepsilon_n}{\varepsilon_{n-1}}
 =\frac{\gamma_n}{\alpha_n}\frac{A_{n-1}}{A_{n+1}}<1;
\]
the strict inequality is the positive $\beta_nA_n$ term in
\eqref{eq:genericrec}.  They tend to zero since the partial sums converge to
$L$.  The alternating-series remainder theorem proves~\eqref{eq:brackets}.
\end{proof}

For the Franel recurrence, Lemma~\ref{lem:cas} gives
\begin{equation}\label{eq:WF}
 F_nC_{n+1}-F_{n+1}C_n=\frac{(-8)^n}{(n+1)^2}.
\end{equation}
Consequently,
\begin{equation}\label{eq:Fseries}
 \boxed{\displaystyle
 \zeta(2)=4\sum_{n=0}^{\infty}
 \frac{(-8)^n}{(n+1)^2F_nF_{n+1}}.}
\end{equation}
This identity, and its connection with a continued fraction discovered by the
Ramanujan Machine, is proved in~\cite{FranelCF}.  Here it also follows at once
from Theorem~\ref{thm:franel} and its Ap\'ery limit.  The additional conclusion
\eqref{eq:brackets} gives certified rational intervals: if
\[
 R_n^{F}=\frac{C_n}{F_n},\qquad
 \varepsilon_n^{F}=\frac{8^n}{(n+1)^2F_nF_{n+1}},
\]
then for even $n$,
$4R_n^F<\zeta(2)<4(R_n^F+\varepsilon_n^F)$, with the odd case given by
the second line of~\eqref{eq:brackets}.

The same calculation for $s_{10}$ appears not to have been recorded in the
sources located for this paper.

\begin{theorem}[$s_{10}$ Casoratian and Basel series]\label{thm:s10cas}
With $S_n,D_n$ as in Theorem~\ref{thm:s10},
\begin{equation}\label{eq:Ws10}
 S_nD_{n+1}-S_{n+1}D_n
 =(-64)^n\frac{(1)_n(3/4)_n(5/4)_n}{(2)_n^3}.
\end{equation}
Hence
\begin{equation}\label{eq:s10series}
 \boxed{\displaystyle
 \zeta(2)=5\sum_{n=0}^{\infty}
 \frac{(-64)^n(1)_n(3/4)_n(5/4)_n}
 {(2)_n^3S_nS_{n+1}}.}
\end{equation}
The partial ratios $D_n/S_n$ give the strict alternating brackets
\eqref{eq:brackets} for $\zeta(2)/5$.
\end{theorem}

\begin{proof}
Here
\[
 \frac{W_n}{W_{n-1}}
 =-\frac{n(64n^2-4)}{(n+1)^3}
 =-64\frac{n(n-1/4)(n+1/4)}{(n+1)^3}.
\]
Iterating from $W_0=1$ proves~\eqref{eq:Ws10}.  Equations
\eqref{eq:s10series} and the brackets follow from Lemma~\ref{lem:cas} and
$D_n/S_n\to\zeta(2)/5$.
\end{proof}

\begin{theorem}[$s_{10}$ polynomial continued fraction]\label{thm:s10cf}
Set
\[
 a_n=(2n+1)(6n^2+6n+2),\qquad
 b_n=n^4(64n^2-4).
\]
Then
\begin{equation}\label{eq:s10cf}
 \boxed{\displaystyle
 \frac5{\zeta(2)}
 =a_0+\cfrac{b_1}{a_1+\cfrac{b_2}{a_2+
             \cfrac{b_3}{a_3+\ddots}}}.}
\end{equation}
The $n$th convergent is exactly $S_{n+1}/D_{n+1}$.
\end{theorem}

\begin{proof}
Let $P_{-2}=0,P_{-1}=1$ and $Q_{-1}=0,Q_0=1$ be the Euler--Wallis numerator
and denominator sequences, so
\[
 P_n=a_nP_{n-1}+b_nP_{n-2}\quad(n\geq0),\qquad
 Q_n=a_nQ_{n-1}+b_nQ_{n-2}\quad(n\geq1).
\]
Multiplying~\eqref{eq:s10rec} by $n!^3$ shows, including the initial cells,
\[
 P_n=(n+1)!^3S_{n+1},\qquad Q_n=(n+1)!^3D_{n+1}.
\]
Thus the $n$th convergent is $P_n/Q_n=S_{n+1}/D_{n+1}$, which tends to
$5/\zeta(2)$.
\end{proof}

The new $s_7$ companion supplies a direct harmonic proof of the continued
fraction evaluation discussed in~\cite{MOCooper}.

\begin{theorem}[$s_7$ Basel series and continued fraction]\label{thm:s7cf}
With $A_n^{(7)},B_n^{(7)}$ as in Theorem~\ref{thm:s7harm},
\begin{equation}\label{eq:Ws7}
 A_n^{(7)}B_{n+1}^{(7)}-A_{n+1}^{(7)}B_n^{(7)}
 =(-27)^n\frac{(1)_n(2/3)_n(4/3)_n}{(2)_n^3}.
\end{equation}
Consequently
\begin{equation}\label{eq:s7series}
 \boxed{\displaystyle
 \zeta(2)=7\sum_{n=0}^{\infty}
 \frac{(-27)^n(1)_n(2/3)_n(4/3)_n}
 {(2)_n^3A_n^{(7)}A_{n+1}^{(7)}}.}
\end{equation}
The partial ratios $B_n^{(7)}/A_n^{(7)}$ give the strict alternating brackets
of Lemma~\ref{lem:cas} for $\zeta(2)/7$.  Moreover, if
\[
 a_n=(2n+1)(13n^2+13n+4),\qquad
 b_n=3n^4(9n^2-1),
\]
then
\begin{equation}\label{eq:s7cf}
 \boxed{\displaystyle
 \frac7{\zeta(2)}
 =4+\cfrac{b_1}{a_1+\cfrac{b_2}{a_2+
              \cfrac{b_3}{a_3+\ddots}}}.}
\end{equation}
Its $n$th convergent is $A_{n+1}^{(7)}/B_{n+1}^{(7)}$.
\end{theorem}

\begin{proof}
Recurrence~\eqref{eq:s7rec} gives
\[
 \frac{W_n}{W_{n-1}}
 =-27\frac{n(n-1/3)(n+1/3)}{(n+1)^3}.
\]
Iteration proves~\eqref{eq:Ws7}; Corollary~\ref{cor:s7limit} and
Lemma~\ref{lem:cas} give the series and brackets.  Multiplication of
\eqref{eq:s7rec} by $n!^3$ gives the Euler--Wallis recurrence with partial
denominators $a_n$ and partial numerators $b_n$.  Its convergents are the
displayed ratios and tend to $7/\zeta(2)$.
\end{proof}

\begin{remark}[arithmetic strength]\label{rem:s7irr}
Theorem~\ref{thm:s7harm} gives a direct harmonic proof of the limit and hence
of the continued-fraction evaluation, but it does not by itself give a new
Ap\'ery proof of the irrationality of $\zeta(2)$.  The characteristic equation
at infinity of~\eqref{eq:s7rec} is
\[
 \lambda^2-26\lambda-27=(\lambda-27)(\lambda+1).
\]
The two formal scales are $27^nn^{-3/2}$ and $(-1)^nn^{-3/2}$.  Thus the
dominant part cancels in
$\zeta(2)A_n^{(7)}-7B_n^{(7)}$, but the remaining linear form is only on the
polynomial scale $O(n^{-3/2})$, rather than exponentially small.  Meanwhile
the evident denominator clearing supplied by~\eqref{eq:s7harm} is
$28\operatorname{lcm}(1,\ldots,2n)^2$.  After that clearing the linear forms
do not tend to zero, so the standard Ap\'ery integrality argument cannot
close.  What the theorem does supply is the harmonic numerator formula, the
companion Lucas law~\eqref{eq:s7lucas}, the exponentially convergent Basel
series~\eqref{eq:s7series}, and the strict rational brackets.
\end{remark}

\section{Comparison with Ap\'ery's \texorpdfstring{$\zeta(3)$}{zeta(3)} companion}

The same program produces a strikingly smaller formula for the numerator row
in Ap\'ery's proof of the irrationality of $\zeta(3)$.  Let
\[
 a_n=\sum_{k=0}^n\bin nk^2\bin{n+k}k^2,
\]
and let $b_0=0,b_1=6$ satisfy Ap\'ery's recurrence
\begin{equation}\label{eq:z3rec}
 (n+2)^3b_{n+2}=(34n^3+153n^2+231n+117)b_{n+1}-(n+1)^3b_n.
\end{equation}
Then the minimal harmonic form is
\begin{equation}\label{eq:z3form}
 b_n=\sum_{k=0}^n\bin nk^2\bin{n+k}k^2
       \left(2H_n^{(3)}-H_k^{(3)}\right).
\end{equation}
An explicit certificate proof and a Lean formalization establish this identity.
Its Lucas descent is the weight-three counterpart of
Theorem~\ref{thm:universal}: multiplication by $p^3$ extracts the high digit of
the harmonic weight.

There is an instructive contrast.  The binomial-power companions use the
symmetric rank-two algebra generated by $1$ and
$\delta=H_k-H_{n-k}$ after differencing.  Formula~\eqref{eq:z3form} is even
smaller: it is linear in order-three harmonic numbers and needs no nested tail.
Neither simplification survives unchanged in the $\zeta(2)$ Ap\'ery pair.

\section{The \texorpdfstring{$\zeta(2)$}{zeta(2)} companion: two closed forms}\label{sec:z2}

Let
\begin{equation}\label{eq:z2kernel}
 T(n,k)=\bin nk^2\bin{n+k}k,\qquad a_n=\sum_{k=0}^nT(n,k),
\end{equation}
and define
\begin{equation}\label{eq:z2op}
 (Lu)_n=(n+1)^2u_{n+1}-(11n^2+11n+3)u_n-n^2u_{n-1}.
\end{equation}
Thus $La=0$.  Let $B_0=0,B_1=1$ and $LB=0$; classically
$B_n/a_n\to\zeta(2)/5$.

\subsection{A proved nested-tail closed form}

Put
\[
 e(n)=\sum_{m=1}^n\frac{(-1)^{m-1}}{m^2},\qquad
 \tau(n,k)=\sum_{m=1}^{\min(k,n)}
 \frac{(-1)^{n+m-1}}{m^2\bin nm\bin{n+m}m}.
\]

\begin{theorem}[$\zeta(2)$ nested-tail formula]\label{thm:z2tail}
For every $n\geq0$,
\begin{equation}\label{eq:z2tail}
 5B_n=\sum_{k=0}^nT(n,k)\bigl(2e(n)+\tau(n,k)\bigr).
\end{equation}
\end{theorem}

The proof is an explicit certificate chain.  Its decisive cancellation can be
stated compactly.  If
\[
 \rho(n,k)=k^2+k(1+6n)-4-15n-11n^2,
\]
then the Zeilberger and Abel residual simplifies to
\begin{equation}\label{eq:miracleW}
 W(n,k)=3-\frac{3(n+1)}{n+1-k}.
\end{equation}
The apparent double pole at $k=n+1$ and the pole at $k=-n$ cancel.  The last
sum is then the elementary binomial identity
\[
 \sum_{k=0}^n\frac{(-1)^k\bin nk}{n+1-k}=\frac{(-1)^n}{n+1}.
\]
The mechanism behind~\eqref{eq:miracleW} is the tail shift
\[
 m^2+(n+1-m)(n+1+m)=(n+1)^2,
\]
which makes the inner $m$-sum telescope.  This proves recurrence membership at
all cells and, with the two initial values, proves
Theorem~\ref{thm:z2tail}.

\subsection{The minimal formula}

The substantially shorter depth-one form is
\begin{equation}\label{eq:z2minimal}
 B_n=\frac15\sum_{k=0}^nT(n,k)
 \left[H_n^{(2)}+H_k\bigl(2H_k-H_{n-k}-H_n\bigr)\right].
\end{equation}
Define
\[
 A_n^\star=\sum_{k=0}^nT(n,k)H_k(2H_k-H_{n-k}-H_n).
\]

\begin{theorem}[minimal $\zeta(2)$ companion]\label{thm:z2minimal}
For every $n\geq0$, formula~\eqref{eq:z2minimal} holds.  Equivalently,
\begin{equation}\label{eq:L2}
 (LA^\star)_n=-a_{n+1}-a_{n-1}\qquad(n\geq1).
\end{equation}
\end{theorem}

\begin{proof}
The harmonic shifts give the exact identity
\[
 L[H_n^{(2)}a_n]=a_{n+1}+a_{n-1}.
\]
Hence applying $L$ to the right side of~\eqref{eq:z2minimal} yields zero
precisely when~\eqref{eq:L2} holds.  The initial values agree.

\medskip
We prove~\eqref{eq:L2} by moving the harmonic letters back into parameters.
Set
\begin{equation}\label{eq:Xmt}
 X_m(t)=\frac{m!}{(1-t)_m}
       =\prod_{j=1}^m\left(1-\frac tj\right)^{-1}
\end{equation}
and introduce the proper hypergeometric sum
\begin{equation}\label{eq:PhiZ2}
 \Phi_n(u,v)=\sum_{k=0}^nT(n,k)X_k(u)
 \frac{X_k(v)^2}{X_{n-k}(v)X_n(v)}.
\end{equation}
At the origin its relevant Taylor coefficients are
\begin{equation}\label{eq:PhiDerivs}
 \Phi_n=a_n,\quad \partial_u\Phi_n=U_n,\quad
 \partial_v\Phi_n=V_n,\quad \partial_u\partial_v\Phi_n=A_n^\star,
\end{equation}
where every derivative is evaluated at $(u,v)=(0,0)$ and
\[
 U_n=\sum_kT(n,k)H_k,\qquad
 V_n=\sum_kT(n,k)(2H_k-H_{n-k}-H_n).
\]

Let $E$ denote the forward shift in $n$ and write
\begin{equation}\label{eq:Lforward}
 \mathcal L=-(n+1)^2-(11n^2+33n+25)E+(n+2)^2E^2.
\end{equation}
Thus $(\mathcal Lu)_n=(Lu)_{n+1}$.  Creative telescoping applied directly to
the summand of~\eqref{eq:PhiZ2} gives an order-three operator
\begin{equation}\label{eq:Cparam}
 \mathcal C(u,v)=\sum_{i=0}^3C_i(n;u,v)E^i
\end{equation}
and a rational certificate $R(n,k;u,v)$ satisfying
\begin{equation}\label{eq:paramCell}
 \sum_{i=0}^3C_i(n;u,v)F(n+i,k;u,v)
 =\Delta_k\bigl(R(n,k;u,v)F(n,k;u,v)\bigr).
\end{equation}
The exact certificate is generated and independently rechecked in
\nolinkurl{work/z2cf/parameter_proof.mac}.  Its denominator in $k$ is
$(n-k+1)^3(n-k+2)^3(n-k+3)^3$.  Summing~\eqref{eq:paramCell} over the natural
finite ranges produces ten top cells and two certificate endpoints.  Their
coefficients of $1,u,v,uv$ reduce respectively to
\begin{equation}\label{eq:paramBoundary}
 0,\qquad0,\qquad0,\qquad0.
\end{equation}

At the origin the operator and its $v$-derivative factor on the right:
\begin{equation}\label{eq:paramFactors}
 \mathcal C(0,0)=Q_0\mathcal L,qquad
 \partial_v\mathcal C(0,0)=Q_v\mathcal L,
\end{equation}
where
\begin{align}\label{eq:Q0Qv}
 Q_0={}&-2(2n+3)(29n+69)+(n+3)(29n+40)E,\\
 Q_v={}&\frac{204n^3+1221n^2+2320n+1383}{(n+1)(n+3)}
 -\frac{138n^2+523n+468}{n+2}E.
\end{align}
The $v$-derivative of the telescoped relation gives $Q_0\mathcal LV=0$,
because $\mathcal La=0$.  Directly $V_0=V_1=V_2=0$, and the leading
coefficient $(n+3)^3(29n+40)$ never vanishes for $n\geq0$; hence
\begin{equation}\label{eq:Vzero}
 V_n=0\qquad(n\geq0).
\end{equation}

Write $C_u=\partial_u\mathcal C(0,0)$ and
$C_{uv}=\partial_u\partial_v\mathcal C(0,0)$.  There are first-order rational
operators $X,Y$ such that
\begin{equation}\label{eq:lclm}
 XQ_v=YQ_0.
\end{equation}
Normalize the coefficient of $E$ in $X$ to $1$; its constant coefficient is
\begin{equation}\label{eq:Xconstant}
 -\frac{2(n+1)(2n+3)
 (10092n^5+133835n^4+700700n^3+1808755n^2+2300030n+1151344)}
 {(n+3)(n+4)
 (10092n^5+83375n^4+266280n^3+408745n^2+299740n+83112)}.
\end{equation}
Exact Ore division gives, for a rational operator $H$,
\begin{equation}\label{eq:decisiveOre}
 YC_u-XC_{uv}-(XQ_0)J=H\mathcal L,qquad J=-(1+E^2).
\end{equation}
The computed right remainder is identically $(0,0,0)$.

Differentiate the telescoped recurrence first in $u$, and then in both $u,v$.
Using~\eqref{eq:Vzero}, one obtains
\[
 Q_0\mathcal LU+C_ua=0,qquad
 Q_0\mathcal LA^\star+Q_v\mathcal LU+C_{uv}a=0.
\]
Multiply these equations by $Y$ and $X$, use~\eqref{eq:lclm}, and subtract.
Equation~\eqref{eq:decisiveOre} and $\mathcal La=0$ then yield
\[
 XQ_0(\mathcal LA^\star-Ja)=0.
\]
The leading coefficient of $XQ_0$ is $(n+4)(29n+69)$, while direct evaluation
at $n=0,1$ gives
$(\mathcal LA^\star-Ja)_0=(\mathcal LA^\star-Ja)_1=0$.  Thus
$\mathcal LA^\star=Ja$, which is precisely~\eqref{eq:L2} after shifting the
index.  This proves Theorem~\ref{thm:z2minimal}.
\end{proof}

\begin{status}[current $\zeta(2)$ status]\label{status:z2}
Both the nested-tail identity~\eqref{eq:z2tail} and the minimal depth-one
identity~\eqref{eq:z2minimal} are now proved by exact symbolic certificates.
The full latter theorem is not yet formalized in Lean.  Its first formal stage,
in \texttt{lean/ZetaLucas/Z2Shell.lean}, contains no \texttt{sorry} and proves
the shell absorption identities, the base Ap\'ery recurrence, and the identity
$L(H_n^{(2)}a_n)=a_{n+1}+a_{n-1}$.  The remaining stage is the certificate and
pre-operator elimination.  Earlier exclusions of three fixed harmonic
certificate modules remain valid: the successful proof escapes them by
telescoping the proper hypergeometric deformation before differentiating.
\end{status}

\section{Structural summary and open problems}

\begin{center}
\small
\begin{tabular}{>{\raggedright\arraybackslash}p{.17\textwidth}
                >{\raggedright\arraybackslash}p{.28\textwidth}
                >{\raggedright\arraybackslash}p{.20\textwidth}
                >{\raggedright\arraybackslash}p{.22\textwidth}}
\toprule
row & closed-form weight & recurrence proof & Lucas/formal status\\
\midrule
Franel $d=3$ & $(\sigma+3\delta^2)/8$ & finite rank-two certificate &
$p^2$ Lucas law; closed form and law in Lean\\
$s_{10}$, $d=4$ & $(\sigma+4\delta^2)/10$ & finite rank-two certificate &
$p^2$ Lucas law; closed form and law in Lean\\
Cooper $s_7$ & equation~\eqref{eq:s7harm} & three polarized parameter
telescopes and Ore elimination & $p^2$ Lucas law; exact symbolic audit;
Lean open\\
Ap\'ery $\zeta(3)$ & $2H_n^{(3)}-H_k^{(3)}$ & finite minimal certificate &
weight-three Lucas descent; Lean proof\\
Ap\'ery $\zeta(2)$ & nested tail and minimal depth-one forms, proved &
finite tail certificate; parameter telescope and Ore elimination &
exact symbolic audit; Lean formalization open\\
\bottomrule
\end{tabular}
\end{center}

The immediate questions are:
\begin{enumerate}
\item Formalize Theorems~\ref{thm:tower}, \ref{thm:s7harm}, and
      \ref{thm:z2minimal} in Lean, including the parameter telescopes and
      Ore-operator identities.
\item Find comparably small all-boundary recurrence certificates for the
      higher deformation rows when $d\geq5$.
\item Strengthen the one-digit congruences to higher powers of $p$ or to a
      compatible multi-digit theory.
\item For $d\geq5$, relate the deformation rows canonically to the
      higher-order recurrence
      solution space rather than merely to the deformation coefficient.
\end{enumerate}

\section*{Reproducibility record}

The formal source is
\texttt{lean/ZetaLucas/FranelClosedForm.lean}; the human-readable formal audit
is \texttt{work/FRANEL\_S10\_LEAN.md}.  The full $\zeta(2)$ tail proof and the
minimal formula are recorded in \texttt{work/z2cf/note.tex}.  The exact
parameter telescope, all boundary checks, operator factorizations, and Ore
remainder are independently recomputed by
\nolinkurl{work/z2cf/parameter_proof.mac}.  For $s_7$, the three exact
telescopes and Taylor certificates are generated by the
\nolinkurl{work/s7_param/extract_*.mac} scripts; the Ore elimination and
right-factor audit are in \nolinkurl{work/s7_param/ore_eliminate.py}.  Running
\nolinkurl{work/s7_param/verify.sh} rebuilds all of them and also performs an
independent exact check of the first fifty rows.  All numerical checks mentioned in
this paper are supplementary evidence only; no theorem labeled ``proved'' uses
finite checking as a proof step.

\begin{thebibliography}{99}

\bibitem{ChamStraub}
M.~Chamberland and A.~Straub,
\emph{Ap\'ery limits: experiments and proofs},
Amer. Math. Monthly \textbf{128} (2021), no.~9, 811--824;
\texttt{arXiv:2011.03400}.

\bibitem{FranelCF}
J.~Tonien,
\emph{Franel numbers and a continued fraction conjecture discovered by the
Ramanujan Machine},
Math. Intelligencer (2026), doi:10.1007/s00283-025-10497-9.

\bibitem{Gessel1982}
I.~M.~Gessel,
\emph{Some congruences for Ap\'ery numbers},
J. Number Theory \textbf{14} (1982), no.~3, 362--368.

\bibitem{Gorodetsky}
O.~Gorodetsky,
\emph{New representations for all sporadic Ap\'ery-like sequences, with
applications to congruences},
Exp. Math. \textbf{32} (2023), no.~4, 641--656;
\texttt{arXiv:2102.11839}.

\bibitem{MalikStraub}
A.~Malik and A.~Straub,
\emph{Divisibility properties of sporadic Ap\'ery-like numbers},
Res. Number Theory \textbf{2} (2016), article 5;
\texttt{arXiv:1508.00297}.

\bibitem{MOCooper}
T.~Piezas III,
\emph{On the continued fractions using Cooper's sequences $s_7,s_{10},s_{18}$
and the Zudilin--Cohen sequence},
MathOverflow question 447166 (2023),
\url{https://mathoverflow.net/questions/447166/}.

\bibitem{Nesterenko1996}
Yu.~V.~Nesterenko,
\emph{A few remarks on $\zeta(3)$},
Math. Notes \textbf{59} (1996), 625--636.

\bibitem{Schneider2007}
C.~Schneider,
\emph{Ap\'ery's double sum is plain sailing indeed},
Electron. J. Combin. \textbf{14} (2007), \#N5.

\bibitem{StraubLucas}
A.~Straub,
\emph{Gessel--Lucas congruences for sporadic sequences},
Monatsh. Math. \textbf{203} (2024), 677--690;
\texttt{arXiv:2301.12248}.

\bibitem{StraubZudilin}
A.~Straub and W.~Zudilin,
\emph{Sums of powers of binomials, their Ap\'ery limits, and Franel's
suspicions},
Int. Math. Res. Not. \textbf{2023}, no.~11, 9861--9879;
\texttt{arXiv:2112.09576}.

\bibitem{vdPoorten1979}
A.~van der Poorten,
\emph{A proof that Euler missed\dots\ Ap\'ery's proof of the irrationality of
$\zeta(3)$},
Math. Intelligencer \textbf{1} (1979), no.~4, 195--203.

\bibitem{Zagier2009}
D.~Zagier,
\emph{Integral solutions of Ap\'ery-like recurrence equations},
in: Groups and Symmetries, CRM Proc. Lecture Notes \textbf{47}, Amer. Math.
Soc., 2009, 349--366.

\end{thebibliography}

\end{document}
